%! Multiplying \& Dividing Radicals; Rationalizing — Unit 6, Lesson 6.3 — Guided Notes (Answer Key)
\documentclass[10pt]{article}
\usepackage{algebra2-article}
\usepackage{algebra2-key}   % pulls in -boxes; do NOT also load -boxes
\newcommand{\TallMath}[1]{$\displaystyle #1\rule[-1.4em]{0pt}{3.2em}$}
% In-formula answer slot (\ans is text-mode and must never appear inside $...$).
\newcommand{\mans}[1]{{\color{keyred}\mathbf{#1}}}
% Vocab answer for the key: bold forest term, then the filled definition on a write-line.
% The leading and trailing \par are required: \ansline ends with \dotfill but does NOT end the
% paragraph, so without them each term label gets pulled onto the previous answer's line.
\newcommand{\vocabans}[2]{%
  \par\noindent\textbf{\textcolor{forest}{#1:}}\\[1pt]\ansline{#2}\par}

\begin{document}

\pageheader{Unit 6, Lesson 6.3}{Guided Notes --- Answer Key}
\namedateperiod

\vspace{0.08in}

\begin{objectivebox}
\textbf{By the end of this lesson, I will be able to\dots}
\begin{itemize}\setlength{\itemsep}{2pt}
  \item \emph{multiply} radical expressions --- including distributing over a sum and multiplying two binomials --- and simplify the result;
  \item recognize a \emph{conjugate} and explain why a conjugate pair always multiplies to a \emph{rational} number;
  \item \emph{rationalize} a denominator, with one radical or with a binomial, and explain why the value of the expression never changes.
\end{itemize}
\end{objectivebox}

\vspace{0.05in}

\begin{vocabbox}
\small
Fill in each term as we build it together.
\par\vspace{2pt}   % \par is required: \vocabans's \noindent is a no-op mid-paragraph
\vocabans{Rationalizing the denominator}{rewriting a fraction so that no radical is left in the denominator, by multiplying the top and bottom by the same expression}
\vocabans{Conjugate}{the same binomial with the middle sign reversed --- the conjugate of $a+\sqrt b$ is $a-\sqrt b$; their product $a^{2}-b$ is always rational}
\vocabans{Multiplying by a form of $1$}{multiplying by a fraction whose numerator and denominator are equal, such as $\sqrt2/\sqrt2$; it changes the \emph{form} of an expression but never its \emph{value}}
\vocabans{Simplest radical form (all three conditions, at last)}{(1) no factor of the radicand is a perfect $n$th power; (2) no fraction under a radical; (3) no radical in a denominator --- condition 3 is today's work}
\end{vocabbox}

\begin{hookbox}
Two students are asked for the exact value of $\dfrac{1}{\sqrt{2}}$. One writes $\dfrac{1}{\sqrt2}$
and stops. The other writes $\dfrac{\sqrt2}{2}$. Their teacher marks \emph{both} correct.
\begin{center}
\renewcommand{\arraystretch}{1.6}
\begin{tabularx}{\linewidth}{@{}X X@{}}
$\dfrac{1}{\sqrt2} \approx \dfrac{1}{1.414\,214}=$ \ans{$0.7071$} &
$\dfrac{\sqrt2}{2} \approx \dfrac{1.414\,214}{2}=$ \ans{$0.7071$} \\
\end{tabularx}
\end{center}
They are the same number. So here is the real question, and it is a fair one:
\textbf{if both are correct, why does every textbook insist on the second one?}
\vspace{2pt}

Imagine it is $1850$ and you must actually \emph{compute} one of these by hand, with no calculator.
Which long division would you rather do --- \; $1 \div 1.414\,214$ \; or \;
$1.414\,214 \div 2$? \ansline{$1.414214\div2$, easily --- halving a number is something you can do
in your head, while dividing by a six-digit decimal is a page of long division. That is the entire
historical reason for the rule: put the ugly number on \emph{top}, where it never has to be divided
by.}
\vspace{2pt}
That is the whole reason the rule exists. A radical on \emph{top} is harmless; a radical on the
\emph{bottom} is a nuisance. Now look at how the second student got there: they multiplied top and
bottom by $\sqrt2$, and $\sqrt2\cdot\sqrt2=$ \ans{$2$} --- the radical simply
\ans{disappeared}. That one fact runs today's entire lesson.
\end{hookbox}

\begin{notesbox}{1. Multiplying: the Product Property, Run Forwards}
Yesterday you read $\sqrt[n]{ab}=\sqrt[n]{a}\cdot\sqrt[n]{b}$ from \emph{left to right} to pull
factors out. Today you read the same statement from \emph{right to left}.
\begin{center}
\begin{tcolorbox}[colback=goldbg, colframe=goldacc, boxrule=0.8pt, arc=1.5mm,
    left=3mm, right=3mm, top=1.5mm, bottom=1.5mm, width=0.92\linewidth]
\centering
$\displaystyle \sqrt[n]{a}\cdot\sqrt[n]{b}=\sqrt[n]{ab}$ \qquad --- \qquad
multiply the \textbf{outsides} together, multiply the \textbf{insides} together, \emph{then}
simplify.
\end{tcolorbox}
\end{center}

\vspace{2pt}
\textbf{Simplify \emph{last}, and never skip it.} A product often contains a perfect power that
neither factor had on its own:
\[
  \sqrt{6}\cdot\sqrt{10}=\sqrt{\mans{60}}=\mans{2\sqrt{15}}
  \qquad\qquad
  \sqrt{3}\cdot\sqrt{12}=\sqrt{\mans{36}}=\mans{6}
\]
Neither $\sqrt6$ nor $\sqrt{10}$ could be simplified alone. Multiplying them created a perfect
square out of nothing --- so \emph{finishing} is not optional.

\vspace{3pt}
\textbf{With coefficients out front.} Outsides multiply outsides; insides multiply insides.
\begin{center}
\renewcommand{\arraystretch}{1.9}
\begin{tabularx}{\linewidth}{@{}X X X@{}}
(a) \; $2\sqrt5\cdot3\sqrt{10}=$ \ans{$30\sqrt2$} & (b) \; $\sqrt[3]{4}\cdot\sqrt[3]{6}=$ \ans{$2\sqrt[3]{3}$} &
(c) \; $\sqrt{2x}\cdot\sqrt{6x^{3}}=$ \ans{$2x^{2}\sqrt3$} \\
\end{tabularx}
\end{center}
{\small\emph{Work:} (a) $6\sqrt{50}=6\cdot5\sqrt2$. \; (b) $\sqrt[3]{24}=\sqrt[3]{8\cdot3}$. \;
(c) $\sqrt{12x^{4}}=\sqrt{4x^{4}}\cdot\sqrt3$.}

\vspace{2pt}
\textbf{The one fact the whole lesson rests on.} A square root multiplied by itself:
\[
  \sqrt{a}\cdot\sqrt{a}=\sqrt{a^{2}}=\mans{a}
  \qquad\text{so}\qquad \left(\sqrt{7}\right)^{2}=\mans{7},
  \qquad \left(3\sqrt5\right)^{2}=9\cdot 5=\mans{45}
\]
\textbf{A square root times itself erases the radical.} (For a cube root you need
\ans{three} copies: $\sqrt[3]{2}\cdot\sqrt[3]{2}\cdot\sqrt[3]{2}=2$. The index still sets the
group size.)

\vspace{3pt}
\emph{One restriction.} The indices must \ans{match} before you may combine two radicals. For
$\sqrt{2}\cdot\sqrt[3]{2}$ there is no radical rule at all --- but Lesson 6.1 handles it:
$2^{1/2}\cdot 2^{1/3}=2^{\,\mans{5/6}}=\sqrt[6]{\mans{32}}$.
\end{notesbox}

\begin{notesbox}{2. Distributing and FOIL --- Radicals Behave Like Binomials}
Nothing new here at all. In Warm-Up item~1 you expanded $3x(2x+5)$ and $(x+4)(x-4)$. Radicals
follow the identical rules; the radical just plays the role of the variable.
\[
  \sqrt{5}\left(\sqrt{3}+\sqrt{10}\right)
  = \sqrt{\mans{15}}+\sqrt{\mans{50}} = \mans{\sqrt{15}+5\sqrt2}
  \qquad
  \sqrt{2}\left(4-\sqrt{2}\right)=\mans{4\sqrt2-2}
\]
\emph{Simplify each product separately after distributing} --- in the first one, only the second
term simplifies.

\vspace{3pt}
\textbf{Two binomials: FOIL, then collect like radicals (Lesson 6.2).}
\[
  \left(2+\sqrt3\right)\left(5-\sqrt3\right)
  = 10 \;\mans{-}\; 2\sqrt3 \;+\; 5\sqrt3 \;\mans{-}\; 3
  = \mans{7+3\sqrt3}
\]

\vspace{3pt}
\textbf{Squaring a binomial --- the error of the day.}
\[
  \left(\sqrt7+\sqrt2\right)^{2}
  = \left(\sqrt7+\sqrt2\right)\left(\sqrt7+\sqrt2\right)
  = 7 + \mans{\sqrt{14}} + \mans{\sqrt{14}} + 2 = \mans{9+2\sqrt{14}}
\]
It is \textbf{not} $7+2=9$. There is a middle term, for exactly the reason
$(x+4)^{2}\ne x^{2}+16$ --- which you already rejected in the Warm-Up. \; Your turn:
\; $\left(3-\sqrt5\right)^{2}=$ \ans{$14-6\sqrt5$}
\end{notesbox}

\begin{notesbox}{3. The Pair That Comes Out Rational: Conjugates}
Now multiply a binomial by its \emph{sign-flipped} partner --- Warm-Up item 1(b), with a radical in
place of the $4$.
\[
  \left(4+\sqrt3\right)\left(4-\sqrt3\right)
  = 16 - 4\sqrt3 + 4\sqrt3 - 3 = \mans{13}
\]
Two things happened at once, and both matter. The middle terms \ans{cancelled}, and the last term
lost its radical because $\sqrt3\cdot\sqrt3=$ \ans{$3$}. What is left is a plain
\ans{rational} number --- no radical anywhere.
\begin{center}
\begin{tcolorbox}[colback=goldbg, colframe=goldacc, boxrule=0.8pt, arc=1.5mm,
    left=3mm, right=3mm, top=1.5mm, bottom=1.5mm, width=0.92\linewidth]
\centering
\textbf{Conjugates:} $a+\sqrt{b}$ and $a-\sqrt{b}$ --- same two terms, \textbf{opposite} middle
sign. \\[3pt]
$\displaystyle \left(a+\sqrt b\right)\left(a-\sqrt b\right)=a^{2}-b$
\qquad
$\displaystyle \left(\sqrt a+\sqrt b\right)\left(\sqrt a-\sqrt b\right)=a-b$ \\[3pt]
{\small Both products are \textbf{always rational}. It is just difference of squares, and squaring
is what undoes a square root.}
\end{tcolorbox}
\end{center}

\vspace{2pt}
\textbf{You have seen this move before.} In Lesson 3.4 you multiplied $a+bi$ by $a-bi$ and got
$a^{2}+b^{2}$ --- a real number. Same trick, different number system: there it turned an
\emph{imaginary} number real; here it turns an \emph{irrational} number rational. \; Why the
different sign? Because $i^{2}=$ \ans{$-1$} while $\left(\sqrt b\right)^{2}=$ \ans{$b$}.

\vspace{3pt}
\textbf{Practice.} Write each conjugate, then multiply the pair.
\begin{center}
\renewcommand{\arraystretch}{1.7}
\begin{tabularx}{\linewidth}{@{}l X X@{}}
\hline
\textbf{Expression} & \textbf{Its conjugate} & \textbf{Their product} \\
\hline
$6+\sqrt{5}$          & \ans{$6-\sqrt{5}$} & \ans{$36-5=31$} \\
$\sqrt{7}-2$          & \ans{$\sqrt{7}+2$} & \ans{$7-4=3$} \\
$\sqrt{11}+\sqrt{3}$  & \ans{$\sqrt{11}-\sqrt{3}$} & \ans{$11-3=8$} \\
\hline
\end{tabularx}
\end{center}
\end{notesbox}

\begin{notesbox}{4. Dividing: the Quotient Property, Run Forwards}
Same idea, one floor down: two radicals with the \emph{same index} can be pulled under one radical.
\begin{center}
\renewcommand{\arraystretch}{2.0}
\begin{tabularx}{\linewidth}{@{}X X X@{}}
(a) \; $\dfrac{\sqrt{48}}{\sqrt{3}}=$ \ans{$4$} &
(b) \; $\dfrac{6\sqrt{20}}{2\sqrt{5}}=$ \ans{$6$} &
(c) \; $\dfrac{\sqrt{40x^{3}}}{\sqrt{5x}}=$ \ans{$2x\sqrt2$} \\
\end{tabularx}
\end{center}
{\small\emph{Work:} (a) $\sqrt{16}$. \; (b) $3\sqrt{4}=3\cdot2$. \;
(c) $\sqrt{8x^{2}}=\sqrt{4x^{2}}\cdot\sqrt2$.}

\vspace{2pt}
Every one of those came out clean. But this one does not:
\; $\dfrac{\sqrt3}{\sqrt5}=\sqrt{\dfrac{3}{5}}$ --- and now there is a radical sitting in the
\ans{denominator}, which breaks simplest-form condition \ans{$3$}. That is the condition we
skipped yesterday. Time to finish it.
\end{notesbox}

\begin{notesbox}{5. Rationalizing a Denominator with One Radical}
\textbf{The move:} multiply the top \emph{and} the bottom by the radical that is stuck in the
bottom. From Warm-Up item~3, that is legal because $\dfrac{\sqrt b}{\sqrt b}=$ \ans{$1$}, and
multiplying by $1$ cannot change a number's value --- only its \ans{form}.
\[
  \frac{1}{\sqrt2}\cdot\frac{\sqrt2}{\sqrt2}
  = \frac{\sqrt2}{\mans{2}}
  \qquad\qquad
  \frac{3}{\sqrt5}\cdot\frac{\mans{\sqrt5}}{\mans{\sqrt5}} = \mans{\frac{3\sqrt5}{5}}
\]
The Hook is now closed: the two students had the same number the whole time.

\vspace{3pt}
\textbf{Your turn.} Rationalize each. In (c), \emph{simplify the denominator first}.
\begin{center}
\renewcommand{\arraystretch}{2.2}
\begin{tabularx}{\linewidth}{@{}X X X@{}}
(a) \; $\dfrac{5}{\sqrt{3}}=$ \ans{$\dfrac{5\sqrt3}{3}$} &
(b) \; $\sqrt{\dfrac{3}{5}}=$ \ans{$\dfrac{\sqrt{15}}{5}$} &
(c) \; $\dfrac{\sqrt{7}}{\sqrt{12}}=$ \ans{$\dfrac{\sqrt{21}}{6}$} \\
\end{tabularx}
\end{center}
{\small\emph{Work:} (b) $\sqrt{3/5}=\dfrac{\sqrt3}{\sqrt5}\cdot\dfrac{\sqrt5}{\sqrt5}
=\dfrac{\sqrt{15}}{5}$. \; (c) $\sqrt{12}=2\sqrt3$ first, then
$\dfrac{\sqrt7}{2\sqrt3}\cdot\dfrac{\sqrt3}{\sqrt3}=\dfrac{\sqrt{21}}{6}$. Multiplying by
$\sqrt{12}$ instead also works, but leaves $\dfrac{\sqrt{84}}{12}$ to clean up.}

\vspace{2pt}
\textbf{A higher index changes what you multiply by.} For $\dfrac{1}{\sqrt[3]{2}}$, multiplying by
$\sqrt[3]{2}$ is \emph{not} enough --- it only gives $\sqrt[3]{4}$ on the bottom, and that is still
a radical. You need \ans{three} equal factors to escape a cube root, so multiply by
$\sqrt[3]{4}$:
\[
  \frac{1}{\sqrt[3]{2}}\cdot\frac{\sqrt[3]{4}}{\sqrt[3]{4}}
  = \frac{\sqrt[3]{4}}{\sqrt[3]{\mans{8}}} = \frac{\sqrt[3]{4}}{\mans{2}}
\]
This is yesterday's \emph{``the index sets the group size''} showing up again: you multiply by
whatever \textbf{completes the group}.
\end{notesbox}

\begin{notesbox}{6. Rationalizing a Binomial Denominator: Use the Conjugate}
If the bottom is $3-\sqrt5$, multiplying by $\sqrt5$ is useless --- try it and you still have a
radical down there. Multiply by the \textbf{conjugate} instead, because a conjugate pair is
guaranteed to come out rational (Section 3).
\[
  \frac{6}{3-\sqrt5}\cdot\frac{\mans{3+\sqrt5}}{\mans{3+\sqrt5}}
  = \frac{6\left(3+\sqrt5\right)}{9-\mans{5}}
  = \frac{6\left(3+\sqrt5\right)}{\mans{4}}
  = \mans{\frac{9+3\sqrt5}{2}}
\]
Notice the last step: the numerator and the denominator shared a common factor, so the answer
reduced. \emph{Always check whether it does.}

\vspace{3pt}
\textbf{Your turn.}
\begin{center}
\renewcommand{\arraystretch}{2.4}
\begin{tabularx}{\linewidth}{@{}X X@{}}
(a) \; $\dfrac{4}{\sqrt5-1}=$ \ans{$\sqrt5+1$} &
(b) \; $\dfrac{\sqrt2}{\sqrt7+\sqrt3}=$ \ans{$\dfrac{\sqrt{14}-\sqrt6}{4}$} \\
\end{tabularx}
\end{center}
{\small\emph{Work:} (a) conjugate $\sqrt5+1$; denominator $5-1=4$, and
$\dfrac{4\left(\sqrt5+1\right)}{4}$ reduces completely. \; (b) conjugate $\sqrt7-\sqrt3$;
denominator $7-3=4$, numerator $\sqrt2\left(\sqrt7-\sqrt3\right)=\sqrt{14}-\sqrt6$.}

\vspace{2pt}
\textbf{Three ways this goes wrong.} Write the fix next to each.
\begin{itemize}[leftmargin=1.5em, itemsep=3pt]
  \item Multiplying only the \emph{denominator} by the conjugate. \; Why that is illegal:
        \ans{that is multiplying by $3+\sqrt5$, not by $1$ --- it changes the value}
  \item Multiplying by the \emph{same} binomial instead of the conjugate. \;
        $\left(3-\sqrt5\right)^{2}=$ \ans{$14-6\sqrt5$}, which is still \ans{irrational}.
  \item Forgetting to FOIL the \emph{numerator} when it is also a binomial. \;
        $\left(2+\sqrt3\right)\left(1+\sqrt3\right)$ has \ans{$4$} products, not $2$.
\end{itemize}
\end{notesbox}

\begin{practicebox}
Work these with your class. Assume all variables represent positive real numbers, and leave no
radical in any denominator.
\begin{enumerate}[leftmargin=1.7em, itemsep=5pt]
  \item Multiply: \; $\sqrt{6}\cdot\sqrt{15}=$ \ans{$3\sqrt{10}$} \quad
        $3\sqrt2\cdot4\sqrt{6}=$ \ans{$24\sqrt3$} \quad $\sqrt[3]{9}\cdot\sqrt[3]{6}=$ \ans{$3\sqrt[3]{2}$}
  \item Expand: \; $\sqrt3\left(\sqrt{12}-\sqrt6\right)=$ \ans{$6-3\sqrt2$} \quad
        $\left(4+\sqrt2\right)\left(3-\sqrt2\right)=$ \ans{$10-\sqrt2$}
  \item Expand: \; $\left(\sqrt5+\sqrt3\right)^{2}=$ \ans{$8+2\sqrt{15}$} \quad
        $\left(2\sqrt7-3\right)\left(2\sqrt7+3\right)=$ \ans{$19$}
  \item Rationalize: \; $\dfrac{8}{\sqrt6}=$ \ans{$\dfrac{4\sqrt6}{3}$} \quad
        $\dfrac{3}{5+\sqrt2}=$ \ans{$\dfrac{15-3\sqrt2}{23}$}
  \item \textbf{Explain.} Both answers in item 3 came from squaring or multiplying a pair of
        binomials with the same two terms, yet one answer still has a radical in it and the other
        does not. What is the difference between the two problems? \ansline{The \textbf{signs}.
        In $\left(\sqrt5+\sqrt3\right)^{2}$ both factors have a plus, so the two middle products
        are both $+\sqrt{15}$ and they \emph{add} to $2\sqrt{15}$ --- the radical survives. In
        $\left(2\sqrt7-3\right)\left(2\sqrt7+3\right)$ the middle signs are opposite, so the middle
        products $+6\sqrt7$ and $-6\sqrt7$ \emph{cancel}, and only the two squares are left, both
        of which are rational. That is exactly what makes the second pair \textbf{conjugates} and
        why conjugates --- not any old pair --- are the tool for clearing a denominator.}
\end{enumerate}
\end{practicebox}

\begin{teachernote}
\textbf{Pacing.} Sections 1--2 are review wearing a new hat and should take about 12 minutes
together; do not linger, because students who can FOIL polynomials can already do this.
\textbf{Section 3 is the intellectual center of the lesson and deserves 10 minutes on its own} ---
the discovery that one specific pair of factors turns an irrational expression rational is the idea
everything after it depends on. Section 5 is quick. Section 6 needs 12 minutes and is where the
lesson is actually assessed.

\textbf{Do not state the conjugate rule before the arithmetic.} Work
$\left(4+\sqrt3\right)\left(4-\sqrt3\right)$ on the board term by term, let the class watch
$-4\sqrt3$ and $+4\sqrt3$ cancel, and only then ask ``what kind of number is $13$?'' The word
\emph{rational} should come from a student. Immediately follow it with
$\left(4+\sqrt3\right)^{2}=19+8\sqrt3$ --- same two numbers, same operation, radical still there.
The contrast is the argument, and it is also the exit ticket's hardest distractor.

\textbf{The Lesson 3.4 callback is worth the ninety seconds.} Students met $(a+bi)(a-bi)=a^{2}+b^{2}$
back in the quadratics unit and used it, if at all, mechanically. Putting it beside
$\left(a+\sqrt b\right)\left(a-\sqrt b\right)=a^{2}-b$ shows that ``multiply by the conjugate'' is a
single general move for getting an unwanted kind of number out of a denominator, not two unrelated
tricks. The sign difference has a reason ($i^{2}=-1$ but $(\sqrt b)^{2}=+b$); make a student say it.

\textbf{Three errors to expect.} (1) \emph{Not simplifying after multiplying} --- $\sqrt6\cdot
\sqrt{10}$ left as $\sqrt{60}$. This is yesterday's ``check the leftover'' habit failing in a new
setting, and it is why the Section 1 examples were chosen so that neither factor simplifies alone.
(2) \emph{Dropping the middle term} in $\left(\sqrt7+\sqrt2\right)^{2}$. Send them straight back to
Warm-Up 1(c); they have already rejected $(x+4)^{2}=x^{2}+16$ today. (3) \emph{Multiplying only the
denominator} by the conjugate --- the tell that a student thinks of rationalizing as a manipulation
of the bottom rather than as multiplying the whole fraction by $1$. That misunderstanding is
diagnosed by the exit ticket's item 2, not by any computation.

\textbf{Section 5's higher-index example is short but load-bearing.} $\dfrac{1}{\sqrt[3]2}$ needs
$\sqrt[3]4$, not $\sqrt[3]2$, and the reason is yesterday's sentence --- the index sets the group
size --- read forwards instead of backwards. Students who got $\sqrt[3]{9}$ wrong on the 6.2 exit
ticket will get this wrong too; they are the same misconception.
\end{teachernote}

\end{document}
